\documentclass[tikz,border=8pt]{standalone}
\usepackage{tikz}
\usepackage{amsmath}
\usepackage{amssymb}
\usepackage{ctex}
\usetikzlibrary{calc, positioning, arrows.meta, decorations.pathreplacing}

\begin{document}
\begin{tikzpicture}[>=Stealth]

%% ===== 左图：线性相关（三个 2D 向量共线）=====
\begin{scope}[xshift=0cm]

  % 背景网格
  \draw[gray!20, thin] (-0.5,-0.5) grid (4.5,3.5);
  \draw[->, gray!60, thin] (-0.5,0) -- (4.8,0) node[right, font=\scriptsize] {$x$};
  \draw[->, gray!60, thin] (0,-0.5) -- (0,3.8) node[above, font=\scriptsize] {$y$};

  % 三个向量共线方向（沿 (2,1) 方向）
  % v1 = (1, 0.5), v2 = (1.5, 0.75), v3 = 2v1 + v2 = (3.5, 1.75) → 截断到合适长度
  % 改用整数/简单分数：v1=(1,1), v2=(2,2)的不同方向，实际共线用(2,1)
  % v1 = (1, 0.5)*2 → 显示 (2,1)
  % v2 = (1, 0.5)*1 → 显示 (1, 0.5)  -- 改为整数: v1=(2,1), v2=(1,0.5)避免小数
  % 用方向(1,0.5)*2, (1,0.5)*3, (1,0.5)*5 -- 但需整数坐标
  % 用方向(2,1): v1=(2,1), v2=(1,0.5) -- 有小数
  % 改用方向(1,1): v1=(1,1), v2=(2,2), v3=3v1=v2+v1=(3,3)？不满足v3=2v1+v2
  % v3=2v1+v2: 若v1=(1,1), v2=(1,1)*alpha, 令v2=(2,2), v3=2*(1,1)+(2,2)=(4,4) 还是共线
  % 标注：v1=(0.8,0.8), v2=(1.2,1.2), v3=2v1+v2=(2.8,2.8) -- 都在y=x线上

  % 用整数坐标：沿(3,2)方向
  % v1 = (1.5,1) -- 避免小数
  % 直接用(3,2)方向: v1=(3,2), v2=(1.5,1), v3=(3,2)+(3,2)+(1.5,1) --有小数
  % 最终: 用(2,1)方向: v1=(2,1), v2=(4,2)=(2*v1), v3=2*v1+v2=2*(2,1)+(4,2)=(8,4) 太大
  % 简单处理: v1=(1,1), v2=(2,2), v3=2*v1+v2=(4,4) --放不下
  % v3=2v1+v2，令v1=(0.6,0.4), v2=(0.9,0.6), v3=2*(0.6,0.4)+(0.9,0.6)=(2.1,1.4)
  % 用(3,2)方向整数倍: v1=(0.3,0.2)*5=(1.5,1)? 还是小数
  % 最终决定: 方向(2,1), v1=(1,0.5)*缩放→整数难，直接用浮点
  % 或：v1=(2,1), v2=(4,2), 那v3=2*v1+v2=(8,4)太大; 用v1=(1,0.5)...
  % 最简：令 v1=(1,1), v2=(1,1), v3=v1+2*v2=(3,3)，标注v3=v1+2v2，都共线y=x

  % 共线向量 along y = x
  \draw[->, very thick, blue!80!black] (0,0) -- (1,1)
    node[above left, font=\small] {$\mathbf{v}_1$};
  \draw[->, very thick, teal!70!black] (0,0) -- (2,2)
    node[above left, font=\small] {$\mathbf{v}_2$};
  \draw[->, very thick, red!70!black] (0,0) -- (3.5,3.5)
    node[below right, font=\small] {$\mathbf{v}_3$};

  % 共线直线提示
  \draw[gray!40, dashed, thin] (-0.3,-0.3) -- (3.8, 3.8);

  % 方程标注
  \node[font=\small, fill=white, draw=gray!40, thin, rounded corners=2pt, inner sep=3pt]
    at (1.0, 3.2) {$\mathbf{v}_3 = \frac{3}{2}\mathbf{v}_2+\frac{1}{2}\mathbf{v}_1$};

  % 原点
  \fill (0,0) circle (1.8pt) node[below left, font=\scriptsize] {$O$};

  % 子图标题
  \node[font=\small\bfseries] at (2.0, -0.9) {线性\textbf{相关}};
  \node[font=\scriptsize, fill=white, draw=gray!50, thin, rounded corners=2pt, inner sep=3pt]
    at (2.0, -1.55) {三向量共线（仅张成 1D 子空间）};

\end{scope}

%% ===== 右图：线性无关（两个非共线 2D 向量）=====
\begin{scope}[xshift=6.5cm]

  % 背景网格
  \draw[gray!20, thin] (-0.5,-0.5) grid (4.5,3.5);
  \draw[->, gray!60, thin] (-0.5,0) -- (4.8,0) node[right, font=\scriptsize] {$x$};
  \draw[->, gray!60, thin] (0,-0.5) -- (0,3.8) node[above, font=\scriptsize] {$y$};

  % 张成平面的网格示意（平行四边形区域）
  \fill[teal!8] (0,0) -- (3,1) -- (4,3) -- (1,2) -- cycle;

  % 两个非共线向量
  \draw[->, very thick, blue!80!black] (0,0) -- (3,1)
    node[below right, font=\small] {$\mathbf{u}_1=(3,1)$};
  \draw[->, very thick, teal!70!black] (0,0) -- (1,2)
    node[above left, font=\small] {$\mathbf{u}_2=(1,2)$};

  % 张成平面提示（在网格范围内的线性组合点）
  \foreach \a/\b in {1/0, 0/1, 1/1}{
    \fill[gray!50] ({3*\a+1*\b}, {1*\a+2*\b}) circle (1.5pt);
  }

  % 原点
  \fill (0,0) circle (1.8pt) node[below left, font=\scriptsize] {$O$};

  % 标注
  \node[font=\small, fill=white, draw=gray!40, thin, rounded corners=2pt, inner sep=3pt]
    at (2.5, 3.2) {张成整个 $\mathbb{R}^2$ 平面};

  % 子图标题
  \node[font=\small\bfseries] at (2.0, -0.9) {线性\textbf{无关}};
  \node[font=\scriptsize, fill=white, draw=gray!50, thin, rounded corners=2pt, inner sep=3pt]
    at (2.0, -1.55) {两向量不共线（张成 2D 平面）};

\end{scope}

%% 整体标题
\node[font=\large\bfseries] at (5.25, 4.8)
  {线性相关 vs 线性无关};

\end{tikzpicture}
\end{document}
